% ============ §1 语音与字母 ============
\hdr{§1}{语音与字母：33 个字母怎么读，怎么连成词}

\callout{\textbf{这一节的坐标}：У1 p9--p30（字母表、元音成对、重音与弱化）、У2 p9--p12（硬软辅音）与 p30（清化表）、У3 p2--p3（重音节律与清浊词表）、p9--p10（调型）、p15--p21（Да／Нет 与 Где／Когда́）。\\
\textbf{这一节要解决的事}：给你一个没见过的俄语词，你能读出它的音、判出它的重音、说清它为什么这么读。}

\subhdr{1.1　33 个字母先分三类}

课件 У1 p9 的字母表把 33 个字母摆成一排，p17 又把元音按「硬—软」成对列出来。两页合起来，字母其实只分三类：

\begin{minipage}{\linewidth}
\tabtop
\begin{tabular}{@{}>{\raggedright\arraybackslash}p{0.14\linewidth}>{\raggedright\arraybackslash}p{0.22\linewidth}>{\raggedright\arraybackslash}p{0.56\linewidth}@{}}
\thead{类别} & \thead{数量} & \thead{内容} \\
\end{tabular}
\tabmid
\begin{tabular}{@{}>{\raggedright\arraybackslash}p{0.14\linewidth}>{\raggedright\arraybackslash}p{0.22\linewidth}>{\raggedright\arraybackslash}p{0.56\linewidth}@{}}
元音字母 & 10（5 对） & а–я ｜ у–ю ｜ э–е ｜ ы–и ｜ о–ё \\
辅音字母 & 21 & б в г д ж з й к л м н п р с т ф х ц ч ш щ \\
无音符号 & 2 & ъ（硬音符号）、ь（软音符号）——本身不发音 \\
\end{tabular}
\tabbot
\end{minipage}

\textbf{这张表在说什么}：它不是三个并列的分类，而是\textbf{一条读音规则的三个零件}。10 个元音刚好成 5 对，每一对的第二个字母都带一个「й 的尾巴」；21 个辅音里有一个 \textbf{й}，它是唯一的「半元音辅音」；两个符号不发音，但它们会改变前面辅音的读法（\S1.4）。先把「а–я 是一对」这件事记住，\S1.2 立刻就用。

\textbf{逐行读法}：第一行 10 个元音，数法是把 \textbf{а я у ю э е ы и о ё} 排开成 5 对，不是从 \textbf{а} 数到第 10 个；第二行 21 个辅音按课件顺序排，\textbf{й} 夹在 \textbf{и} 后面（它长得像 \textbf{и}，读音也像 [й]）；第三行两个符号永远不出现在词首。

\aside{课件 У1 p9 只给了字母表面，没有给字母名称（бэ、вэ、эль……）。本讲义按标准俄语字母名称补一张读音表：元音字母名＝它本身的音（а [а]、о [о]、у [у]、ы [ы]、э [э]），辅音字母名＝辅音＋э（бэ、вэ、гэ、дэ、жэ、зэ、пэ、тэ、фэ、цэ、че、ша、ща、ка、эль、эм、эн、эр、эс、ха）；й 叫「и кра́ткое」（短 и），ь 叫「мя́гкий знак」（软音符），ъ 叫「твёрдый знак」（硬音符），я 叫 [йа]、ю 叫 [йу]、е 叫 [йэ]、ё 叫 [йо]。字母名称只用于拼读单词（如「эс-тэ-о-эль」拼 стол），不进入句子。}

\subhdr{1.2　元音 5 对：每一对的差别，是「有没有 й」}

У1 p17 把五对元音并排列出：\textbf{а–я、у–ю、э–е、ы–и、о–ё}。课件只给了这张对照，没说这一对差在哪。差的是同一个东西：

\begin{minipage}{\linewidth}
\tabtop
\begin{tabular}{@{}>{\raggedright\arraybackslash}p{0.11\linewidth}>{\raggedright\arraybackslash}p{0.17\linewidth}>{\raggedright\arraybackslash}p{0.19\linewidth}>{\raggedright\arraybackslash}p{0.44\linewidth}@{}}
\thead{硬} & \thead{软} & \thead{软音的来源} & \thead{在词首或元音后怎么读} \\
\end{tabular}
\tabmid
\begin{tabular}{@{}>{\raggedright\arraybackslash}p{0.11\linewidth}>{\raggedright\arraybackslash}p{0.17\linewidth}>{\raggedright\arraybackslash}p{0.19\linewidth}>{\raggedright\arraybackslash}p{0.44\linewidth}@{}}
а & я & [й] ＋ [а] & я 读 [йа]：я́блоко 的第一个音是 [йа] \\
у & ю & [й] ＋ [у] & ю 读 [йу]：ю́г 读 [йук] \\
э & е & [й] ＋ [э] & е 读 [йэ]：е́сли 的第一个音是 [йэ] \\
ы & и & 无 [й]，纯元音 & 这一对是例外：и 不是 [йы]，它就是 [и] \\
о & ё & [й] ＋ [о] & ё 读 [йо]：ёлка 读 [йолка] \\
\end{tabular}
\tabbot
\end{minipage}

\textbf{这张表在说什么}：它把「五对元音」翻译成一句话——\textbf{я、ю、е、ё 这四个字母，本身就是「й ＋ 元音」}；\textbf{ы／и} 是唯一不带 й 的一对，它们的差别在舌位，不在有没有 й。

\textbf{逐行读法}：前三行（я／ю／е）和第五行（ё）的结构完全一样：[й] 起头，后面接一个元音。第四行（ы／и）必须单独记：\textbf{ы} 是俄语特有的音，汉语里没有，读法是发 [и] 时把舌位往后压。

\textbf{关键的一步在下一节}：这四个「带 й 的元音」一旦跑到\textbf{辅音后面}，那个 [й] 就不读了，取而代之的是「把前面那个辅音读软」。所以 \textbf{тя} 不是 [тйа]，而是 [т'а]。这就是 \S1.4 的全部内容。

\warnbox{\textbf{边界}：\textbf{и} 虽然是「ы 的软对」，但它并不带 [й]。常见错法是把它也当 [йи] 去读——不对，\textbf{кни́га} 读 [кн'ига]，中间没有 [й]。}

\subhdr{1.3　重音：一个词只有一个，且它决定了怎么读}

У1 p28--p30 把 \textbf{о} 分成两组：一组读饱满的 [о]，一组弱化成 [а]。分组的标准不是词性、不是位置，而是\textbf{重音在不在它上面}：

\begin{minipage}{\linewidth}
\tabtop
\begin{tabular}{@{}>{\raggedright\arraybackslash}p{0.24\linewidth}>{\raggedright\arraybackslash}p{0.20\linewidth}>{\raggedright\arraybackslash}p{0.44\linewidth}@{}}
\thead{情况} & \thead{怎么读} & \thead{课件例词} \\
\end{tabular}
\tabmid
\begin{tabular}{@{}>{\raggedright\arraybackslash}p{0.24\linewidth}>{\raggedright\arraybackslash}p{0.20\linewidth}>{\raggedright\arraybackslash}p{0.44\linewidth}@{}}
о 带重音 & 读饱满的 [о] & до́м, кот, сто́л, о́н, мно́го, фо́то, до́ма \\
о 不带重音 & 弱化成 [а] & вода́, когда́, пока́, нога́, она́, хорошо́ \\
整个词 & 只有一个音节重读 & молоко́（重音落在最后一个 о） \\
\end{tabular}
\tabbot
\end{minipage}

\textbf{这张表在说什么}：它只讲了一件事——\textbf{同一个字母 о，重读时读 [о]，不重读时读 [а]}。因此「怎么读一个 о」这个问题，不能看着字母回答，必须先找到重音。

\textbf{逐行读法}：第一行七个词，重音都在 о 上，所以读 [о]；第二行六个词，重音都不在第一个 о 上（вода́ 重音在 а、хорошо́ 重音在最后的 о），那个不重读的 о 读 [а]。第三行 \textbf{молоко́} 是这一页最能说明问题的一个词：三个 о，只有最后一个是 [о]，前两个都弱化成 [а]，读成 [малако́]。

\textbf{补：为什么会弱化}。课件只给词表，没给原因。原因是\textbf{重读音节的气流最强}，非重读的元音发不满，口型自然向 [а] 靠（语言学术语叫「元音弱化」，реду́кция）。这条不必背术语，只要记住操作顺序：\textbf{先找重音 → 再看 о 在不在重音上 → 不在就念 [а]}。

\aside{这条规则是 \S1.5 清化的「邻居」：一个管元音（о 变 [а]），一个管辅音（浊变清）。两者共用同一个前提——重音位置。所以本讲义从 \S1.3 起，所有俄语词都标重音。}

\subhdr{1.4　硬辅音与软辅音：软，是「被后面的元音带软的」}

У2 p10--p12 用三列一组的对照做练习：\textbf{та — тя — тья}、\textbf{да — дя — дья}、\textbf{по — пё — пьё}、\textbf{лэ — ле — лье}。三列的关系不是三套词尾，而是\textbf{同一个辅音的三种硬度}：

\begin{minipage}{\linewidth}
\tabtop
\begin{tabular}{@{}>{\raggedright\arraybackslash}p{0.12\linewidth}>{\raggedright\arraybackslash}p{0.26\linewidth}>{\raggedright\arraybackslash}p{0.26\linewidth}>{\raggedright\arraybackslash}p{0.26\linewidth}@{}}
\thead{组} & \thead{硬} & \thead{软（元音带的）} & \thead{更软（加 ь）} \\
\end{tabular}
\tabmid
\begin{tabular}{@{}>{\raggedright\arraybackslash}p{0.12\linewidth}>{\raggedright\arraybackslash}p{0.26\linewidth}>{\raggedright\arraybackslash}p{0.26\linewidth}>{\raggedright\arraybackslash}p{0.26\linewidth}@{}}
т 组 & та & тя & тья \\
д 组 & да & дя & дья \\
п 组 & по & пё & пьё \\
л 组 & лэ & ле & лье \\
\end{tabular}
\tabbot
\end{minipage}

\textbf{这张表在说什么}：三列的差别只有一个动作——\textbf{舌面抬不抬起来}。抬起来（第二、三列）就是软辅音，不抬（第一列）就是硬辅音。抬的动作由谁发起？由\textbf{后面的元音}发起：\textbf{я、е、ё、ю、и} 这五个元音，以及软音符 \textbf{ь}，都会让前面的辅音变软。

\textbf{逐行读法}：第一列 \textbf{та} 读 [та]，辅音是硬的；第二列 \textbf{тя} 读 [т'а]，这里的 я 不再读 [йа]，它的 [й] 掉了，只留下「把 т 读软」这一个作用；第三列 \textbf{тья} 读 [т'йа]，软音符 ь 让 т 先软一次，я 又把它自己的 [й] 读出来，所以出现了两个音。

\textbf{补：\textbf{ш} 和 \textbf{ж} 是唯一不听话的两个}。У2 p12 把它们单独列出来：\textbf{ши} 读 [шы]、\textbf{жи} 读 [жы]、\textbf{ше} 读 [шэ]、\textbf{же} 读 [жэ]。原因是这两个音在俄语里\textbf{永远硬}，没有软的对立面，所以 и 到了它们后面只能读 [ы]，е 到了它们后面只能读 [э]。这条规则在 \S4.2 讲形容词拼写时还会再出现一次（ж／ш／ч／щ 后面写 и 不写 ы）。

\warnbox{\textbf{边界一}：不是所有辅音都有软硬两版。\textbf{ш、ж、ц} 永远硬；\textbf{ч、щ、й} 永远软。\\
\textbf{边界二}：软音不是「加一个 [ь] 音」。\textbf{ть} 和 \textbf{т} 的差别不在多了一个音，而在辅音本身的音色——这一点用汉语拼音的送气／不送气类比更容易体会：音还是那个音，只是舌位不同。}

\subhdr{1.5　清浊六对与两条清化：一条规律的两种位置}

У3 p2--p3 给的是六对清浊辅音和词末清化的例词；У2 p30 给的是「浊辅音站在清辅音前面」的九组组合。两份课件隔了一讲，其实说的是同一条规律：

\begin{minipage}{\linewidth}
\tabtop
\begin{tabular}{@{}>{\raggedright\arraybackslash}p{0.14\linewidth}>{\raggedright\arraybackslash}p{0.14\linewidth}>{\raggedright\arraybackslash}p{0.14\linewidth}>{\raggedright\arraybackslash}p{0.14\linewidth}>{\raggedright\arraybackslash}p{0.14\linewidth}>{\raggedright\arraybackslash}p{0.14\linewidth}@{}}
\thead{浊} & \thead{清} & \thead{浊} & \thead{清} & \thead{浊} & \thead{清} \\
\end{tabular}
\tabmid
\begin{tabular}{@{}>{\raggedright\arraybackslash}p{0.14\linewidth}>{\raggedright\arraybackslash}p{0.14\linewidth}>{\raggedright\arraybackslash}p{0.14\linewidth}>{\raggedright\arraybackslash}p{0.14\linewidth}>{\raggedright\arraybackslash}p{0.14\linewidth}>{\raggedright\arraybackslash}p{0.14\linewidth}@{}}
б & п & в & ф & г & к \\
д & т & ж & ш & з & с \\
\end{tabular}
\tabbot
\end{minipage}

\textbf{这张表在说什么}：六对上下对齐——把上面那个读成下面那个，就是把浊辅音「清化」。清化的触发条件有两个，课件分在两处讲，本讲义合并成一条：

\begin{itemize}
\item \textbf{位置一：词末}（У3 p2--p3）。\textbf{друг} 读 [друк]、\textbf{клуб} 读 [клуп]、\textbf{эта́ж} 读 [эта́ш]——г→[к]、б→[п]、ж→[ш]。
\item \textbf{位置二：浊辅音站在清辅音前面}（У2 p30）。九组：\textbf{бк}→[пк]（улы́бка）、\textbf{бс}→[пс]、\textbf{вт}→[фт]（вто́рник、за́втра）、\textbf{вк}→[фк]（остано́вка）、\textbf{вс}→[фс]（все）、\textbf{вх}→[фх]（вход）、\textbf{вч}→[фч]（вчера́）、\textbf{дс}→[тс]、\textbf{жк}→[шк]（ло́жка）。
\end{itemize}

\textbf{逐行读法}：九组里出现最多的两个浊辅音是 \textbf{б} 和 \textbf{в}，因为俄语里最常见的词尾就是 \textbf{-бка}、\textbf{-вка}、\textbf{-втор}。判断顺序固定：\textbf{先找这个浊辅音右边有没有清辅音或词界，有就按表清化，没有就读本音}。

\textbf{补：为什么只有这六个}。清浊成对的只有 б／п、в／ф、г／к、д／т、ж／ш、з／с 六组；\textbf{л、м、н、р、й} 是响音，声带一直在振动，不参与清化，所以 \textbf{стол} 里的 л 照读、\textbf{парк} 里的 р 照读。这条是 \S1.5 表里为什么只有六行的原因。

\warnbox{\textbf{边界}：в 在 \textbf{вт、вс、вх、вч、вк} 里清化，但在 \textbf{два}、\textbf{твой} 这类词里，它后面跟的是浊辅音或响音，就照读 [в]——不是「в 见谁都变 [ф]」。}

\subhdr{1.6　三个调型：同一句话，调子不同，意思不同}

У3 p9--p10 给了 ИК-1 与 ИК-3，У8 p27 补了 ИК-5。三个调型对应三种句子功能，共同点是：\textbf{调型的「中心」落在你要强调的那一个词上}。

\begin{minipage}{\linewidth}
\tabtop
\begin{tabular}{@{}>{\raggedright\arraybackslash}p{0.10\linewidth}>{\raggedright\arraybackslash}p{0.24\linewidth}>{\raggedright\arraybackslash}p{0.28\linewidth}>{\raggedright\arraybackslash}p{0.30\linewidth}@{}}
\thead{调型} & \thead{用在哪} & \thead{声音怎么走} & \thead{例} \\
\end{tabular}
\tabmid
\begin{tabular}{@{}>{\raggedright\arraybackslash}p{0.10\linewidth}>{\raggedright\arraybackslash}p{0.24\linewidth}>{\raggedright\arraybackslash}p{0.28\linewidth}>{\raggedright\arraybackslash}p{0.30\linewidth}@{}}
ИК-1 & 陈述句 & 中心词上平降，句末落到底 & Э́то А́нна. ｜ Он до́ма. \\
ИК-3 & 没有疑问词的一般疑问句 & 中心词的重读元音急升，句尾再降 & Это Анто́н? \\
ИК-5 & 感叹句 & 中心词先升后降，音域拉大 & Како́е си́нее мо́ре! \\
\end{tabular}
\tabbot
\end{minipage}

\textbf{这张表在说什么}：三个调型的差别只在「音高怎么走」，词一个都不换。\textbf{Э́то Ива́н.}（ИК-1）是陈述，\textbf{Э́то Ива́н?}（ИК-3）是提问——写出来只差一个标点，读出来的差别全在调子上。

\textbf{逐行读法}：ИК-1 的落点在句末，所以你听到声音往下沉，就知道话说完了；ИК-3 的升调不在句末，而在\textbf{被问的那一项}上，所以 \textbf{Это ³Анна?}（问「是 Анна 吗」）和 \textbf{³Это Анна?}（问「这是 Анна 吗」）听起来是不同的两个问题；ИК-5 的音域比前两个都大，句尾带感叹。

\aside{课件的标注只画了中心位置（У3 p10 在句子上打「³」），没有说明「中心移动＝意思移动」。这一层是本讲义补的：ИК-3 的中心永远落在\emph{你要问的那一项}上，中心一挪，问的东西就变了。}

\warnbox{\textbf{边界}：有疑问词的问句（\textbf{Кто э́то?}、\textbf{Где банк?}）不用 ИК-3，用 ИК-2（课件未讲，本讲义不展开）。所以「问句都用上扬调」这句话是错的。}

\subhdr{1.7　五个基本问答：这一节就把它们说完}

前六节全是零件，这一节把它们装成句子。课件里反复操练的就是下面五组（У2 p19--p25、У3 p15--p21、У4 p9）：

\begin{minipage}{\linewidth}
\tabtop
\begin{tabular}{@{}>{\raggedright\arraybackslash}p{0.30\linewidth}>{\raggedright\arraybackslash}p{0.34\linewidth}>{\raggedright\arraybackslash}p{0.26\linewidth}@{}}
\thead{问} & \thead{问的是什么} & \thead{答} \\
\end{tabular}
\tabmid
\begin{tabular}{@{}>{\raggedright\arraybackslash}p{0.30\linewidth}>{\raggedright\arraybackslash}p{0.34\linewidth}>{\raggedright\arraybackslash}p{0.26\linewidth}@{}}
Кто э́то? & 人、动物 & Э́то студе́нт. ｜ Э́то соба́ка. \\
Что э́то? & 东西、事情 & Э́то фру́кты. ｜ Э́то кни́га. \\
Э́то …? & 是不是（一般疑问） & Да, э́то … ｜ Нет, э́то не … \\
Где …? & 地点 & до́ма ｜ тут ｜ там ｜ сле́ва ｜ спра́ва ｜ впереди́ \\
Когда́ …? & 时间 & у́тром ｜ днём ｜ ве́чером ｜ но́чью ｜ зимо́й ｜ ле́том \\
\end{tabular}
\tabbot
\end{minipage}

\textbf{这张表在说什么}：五组问答共用同一个骨架——\textbf{疑问词 ＋ э́то}，回答一律用 \textbf{Э́то ＋ 名词}。第二组（Что э́то?）与第一组（Кто э́то?）的分界线只有一个：\textbf{问的是人／动物，还是东西}。

\textbf{逐行读法}：第三行是一般疑问，问句和陈述句写法一样，只改调型（\S1.6）；第四、五行答的不是名词而是副词，这类副词在句子里\textbf{不变形}——\textbf{Где банк? — Банк сле́ва.} 里的 сле́ва 永远是 сле́ва，不随 банк 的性数变。

\textbf{补充：否定的两种写法}。У3 p15--p16 给了两条路：只换词（\textbf{Нет, э́то чай.}），或先否定原词再给新词（\textbf{Нет, э́то не студе́нт. Э́то преподава́тель.}）。后一种信息更全，正式场合优先。

\subhdr{1.8　练习：六题，答案是「怎么想出来的」}

课件 У1--У3 的练习以朗读为主（У2 p10--p12 的音节对照、У3 p3 的词表）。下面六题把那些朗读练习转成可以落笔的题，答案依据 \S1.1--\S1.7 的规则。

\begin{enumerate}[label=\textbf{第 \arabic* 题}]
\item （У2 p30 原表）注音：ло́жка ｜ все ｜ вто́рник
\item （У2 p10--p12）下面哪一组是软辅音？та–ту–ты ｜ тя–тю–ти
\item （У3 p2--p3）给清辅音配出浊辅音：п–\blank ｜ т–\blank ｜ с–\blank
\item （У3 p9--p10、У8 p27）判断调型：Э́то тетра́дь. ｜ Э́то тетра́дь? ｜ Како́й краси́вый го́род!
\item （У2 p19--p25）用 Кто 或 Что 填空：— \blank э́то? — Э́то преподава́тель. ｜ — \blank э́то? — Э́то молоко́.
\item （У3 p17--p21）补副词：— Где банк? — Банк \blank（右边）。 ｜ — Когда́ обе́д? — Обе́д \blank（白天）。
\end{enumerate}

\textbf{第 1 题}　关键一步：先看「浊辅音右边有没有清辅音」。ж 右边是 к、в 右边是 с、в 右边是 т，三处都要清化。答案：ло́жка [ло́шка] ｜ все [фсэ] ｜ вто́рник [фто́рник]。

\textbf{第 2 题}　关键一步：看元音是不是 я／ю／и 这三个「带软」的字母。答案：тя–тю–ти——三个元音都把前面的 т 读软成 [т']。

\textbf{第 3 题}　关键一步：按 \S1.5 的六对表横着找。答案：п–б、т–д、с–з。

\textbf{第 4 题}　关键一步：看标点和疑问词——句号、无疑问词的问号、感叹号，分别对应 ИК-1、ИК-3、ИК-5。答案：ИК-1 ｜ ИК-3 ｜ ИК-5。

\textbf{第 5 题}　关键一步：先判答句里那个名词是「人」还是「东西」。答案：Кто（преподава́тель 是人）｜ Что（молоко́ 是东西）。

\textbf{第 6 题}　关键一步：先看疑问词是 Где（地点）还是 Когда́（时间），再去对应那一列找词。答案：спра́ва ｜ днём。
